\chapter{Используемое в книге программное обеспечение}

\section{SageMath}
\label{sagemath_setup}
SageMath, также известная как Sage, — это свободная программная система (software system) с открытым исходным кодом (open-source), предоставляющая обширный набор математических инструментов и возможностей. Она предлагает унифицированный интерфейс к различным математическим библиотекам (libraries) и инструментам, включая арифметику конечных полей, эллиптические кривые и криптографические примитивы, среди прочего. С помощью SageMath пользователи получают доступ к мощной и эффективной платформе для выполнения вычислений, анализа данных и визуализации результатов в широком спектре математических областей. Чтобы обеспечить интерактивный опыт обучения и дать возможность поработать руками с концепциями, описанными в этой книге, мы приводим примеры того, как программировать их в \href{https://www.sagemath.org/}{Sage}. Sage основан на удобном для обучения языке программирования \href{https://www.python.org/}{Python}, расширенном и оптимизированном для вычислений с алгебраическими объектами. Поэтому мы рекомендуем выполнить установку (installation) Sage, прежде чем погружаться в следующие главы.

Шаги установки для различных конфигураций системы описаны на \href{https://doc.sagemath.org/html/en/installation/index.html}{сайте Sage}. Обратите внимание, что мы используем версию Sage 9, поэтому если вы используете Linux и ваш менеджер пакетов содержит только версию 8, вам может потребоваться выбрать другой путь установки, например, использование предварительно собранных бинарных файлов. Если вы не знакомы с SageMath, мы рекомендуем вам обратиться к \href{https://doc.sagemath.org/html/en/tutorial/index.html}{руководству Sage}.

\section{Circom}
\label{circom_setup}
Circom — это язык программирования и компилятор (compiler) для проектирования арифметических схем (circuits). Он предоставляет платформу для программистов, чтобы создавать собственные схемы (circuits), которые затем могут быть скомпилированы в системы ограничений ранга 1 (Rank-1 Constraint Systems, R1CS) и выведены как программы WebAssembly и C++ для эффективного вычисления (efficient evaluation). Библиотека с открытым исходным кодом CIRCOMLIB предлагает набор готовых шаблонов схем (circuit templates), которые могут быть использованы пользователями Circom.

Чтобы скомпилировать наши вычисления на бумаге в реальные zk-SNARK, мы приводим примеры того, как реализовать их с помощью Circom. Поэтому рекомендуется установить Circom перед изучением следующих глав. Для генерации и проверки zk-SNARK для наших схем Circom мы используем библиотеку SNARK.js, которая является реализацией схем ZK-SNARK на JavaScript и чистом WebAssembly. Она использует протокол Groth16 и PLONK. Шаги установки можно найти на странице установки Circom \href{https://docs.circom.io/getting-started/installation/#installing-circom}{Circom installation}.
